\subsection{The first obstruction beyond the splitting of \S51.12}

\begin{statement}{E.26}{Proved}{The class-two splitting obstruction is the generator modulo three}
Let
\[
 P=\langle\rho,\sigma\mid\rho^3=\sigma^2=1\rangle
   =\operatorname{PSL}_2(\mathbb Z),\qquad
 q:P\longrightarrow S_3,
\]
where $q(\rho)$ is a three-cycle and $q(\sigma)$ a transposition.
The kernel is the free group
\[
 F=\langle x,y\rangle,\qquad
 x=(\sigma\rho)^2,\quad y=(\rho\sigma)^2.
\]
The extension of $S_3$ by $F/F'$ splits by B.11. Its class-two
refinement
\[
 1\longrightarrow F/\gamma_3F\longrightarrow P/\gamma_3F
       \longrightarrow S_3\longrightarrow1                         \tag{E.26a}
\]
does not split. Neither does its profinite completion, its relative
pro-$3$ completion, or any of the discrete or profinite extensions
$P/F^{(j)}\to S_3$ for $j\ge2$.

Write $c=[x,y]$ and $r'=x\rho$. The inverse image in (E.26a) of
the chosen abelian complement is
\[
 1\longrightarrow\mathbb Z(\mathrm{sign})
   \xrightarrow{\,1\mapsto c\,}
   \langle r',\sigma\mid\sigma^2=1,\ \sigma r'\sigma=(r')^{-1}\rangle
   \longrightarrow S_3\longrightarrow1,
 \qquad (r')^3=c.                                             \tag{E.26b}
\]
Its cohomology class is $1$ in
$H^2(S_3,\mathbb Z(\mathrm{sign}))\simeq\mathbb Z/3$.
The finite quotient
\[
 P/(\gamma_3F\,F^3)\longrightarrow S_3                          \tag{E.26c}
\]
has order $162$, Heisenberg kernel of order $27$, and no section.
The inverse image of the chosen abelian complement is the dihedral
group of order $18$.

For the class-two relative pro-$p$ extension, the set of continuous
sections modulo conjugation by its kernel is empty at $p=3$, has
two elements at $p=2$, and has one element at every $p\ge5$.
Over $\mathbb Q$ or $\mathbb Q_\ell$, the rational class-two
relative extension does split, with one conjugacy class of sections.
An explicit section in every one of these cases where it exists is
\[
 \bar\rho\longmapsto c^{-1/3}x\rho,
 \qquad \bar\sigma\longmapsto\sigma.                          \tag{E.26d}
\]
Here $c^a$ denotes the actual central coordinate $a$ of the
class-two kernel; it is defined over $\mathbb Z_p$ for $p\ne3$
and over the indicated fields.

\emph{Manuscript: \S51.12, folios 700--704 and 715--716,
scans 105--114 and 135--138; in particular the abelian splitting
on scan 136. This answers the splitting part of E(c) with its
integral, profinite, and rational comparison maps.}
\end{statement}
\begin{proof}
The free-kernel identification and the matrices of its conjugation
action were proved in A.2 and B.9. In the present central quotient
of the modular group the exact free-group substitutions are
\[
 \rho x\rho^{-1}=y,\qquad
 \rho y\rho^{-1}=x^{-1}y^{-1},\qquad
 \sigma x\sigma=y,\quad\sigma y\sigma=x.                      \tag{E.26e}
\]
They also follow immediately by substituting the displayed words
for $x,y$ and using $\rho^3=\sigma^2=1$.
Consequently, already in $P$,
\[
 (x\rho)^3=xyx^{-1}y^{-1}=c,
 \qquad ((x\rho)\sigma)^2=1.                                 \tag{E.26f}
\]
The second identity follows alternatively by reducing
$(\sigma\rho\sigma\rho^2\sigma)^2$.

We specify the integral class-two group rather than replacing it
by a vector space. Every element of $F/\gamma_3F$ has a unique
normal form $x^a y^b c^k$, $a,b,k\in\mathbb Z$, with law
\[
 (a,b,k)(a',b',k')
   =(a+a',b+b',k+k'-a'b).                                    \tag{E.26g}
\]
Existence is collection using the central commutator $c$.
Uniqueness follows from the injective matrix model
\[
 x^a y^b c^k\longmapsto
 \begin{pmatrix}1&a&ab+k\\0&1&b\\0&0&1\end{pmatrix}.
\]
Thus $c$ has infinite order. Conjugation by $r'$ fixes $c$ and
conjugation by $\sigma$ sends $c$ to $c^{-1}$. The inverse image
of the complement generated by $\overline{x\rho},\bar\sigma$
in $P/F'$ has elements $c^k(r')^i\sigma^j$, with
$0\le i<3$, $0\le j<2$. Since $c=(r')^3$, these are exactly
the normal forms of the infinite dihedral group in (E.26b).
There are no additional relations: the image in $S_3$ first
determines $i,j$, and $c^k=1$ then forces $k=0$.

For clarity we compute the cohomology class. Write elements of
$S_3$ as $(i,j)=\bar\rho^i\bar\sigma^j$, with $0\le i<3$ and
$0\le j<2$. The section $(i,j)\mapsto(r')^i\sigma^j$ gives the
normalized integral two-cocycle
\[
 f((i,j),(i',j'))
   =\left\lfloor\frac{i+(-1)^j i'}{3}\right\rfloor,
 \qquad (i,j)\cdot a=(-1)^j a.                               \tag{E.26h}
\]
Its restriction to the cyclic subgroup has carry $1$, namely
$(r')^3=c$. We can compute the whole cohomology group directly.
In an arbitrary extension by $A(\mathrm{sign})$, for
$A=\mathbb Z$ or $\mathbb Z_3$, choose lifts $r,s$. The elements
$s^2$ and $(rs)^2$ lie in $A$ and are fixed by conjugation by
$s$ and $rs$, respectively. Those conjugations act as $-1$ on
$A$; because $A$ has no $2$-torsion, both elements vanish.
The only remaining parameter is $r^3=a\in A$. Replacing $r$
by $b r$, with additive notation on $A$, changes $a$ to $a+3b$.
Conversely every $a\in A$ defines an extension: take
$(A\oplus\mathbb Zr)/(3r-a)$ and adjoin the involution acting
by negation. This proves $H^2(S_3,A(\mathrm{sign}))=A/3A$ and
identifies (E.26b) with $1$. Equivalently,
the one-cochain $(i,j)\mapsto i$ has coboundary $3f$.

To exclude sections of the whole group, we must also exclude
other abelian complements. B.10 proves, over every coefficient
group used here,
\[
 H^0(S_3,M)=H^1(S_3,M)=0,
 \qquad M=F^{\mathrm{ab}}.
\]
Thus conjugation by a unique element of $M$ carries any abelian
section to the B.11 section. The surjection
$F/\gamma_3F\to M$ lifts this conjugating element.
After this conjugation any proposed section of (E.26a) must send
$\bar\rho$ to $c^a r'$ and $\bar\sigma$ to $c^b\sigma$.
Its order-three relation is
\[
 (c^a r')^3=c^{1+3a}=1.                                    \tag{E.26i}
\]
There is no solution $a\in\mathbb Z$ or $a\in\mathbb Z_3$.
The other two defining relations hold for all $a,b$, because
$\sigma$ and $r'\sigma$ invert $c$ and have order two.
This proves the integral and relative pro-$3$ assertions.

Here are the exact completion and quotient comparisons. Since
$F$ is normal in $P$ and $\gamma_3F$, $F^3$, and $F^{(j)}$ are
characteristic in $F$, every displayed quotient is a group
extension over the same $S_3$. The relative pro-$p$ class-two
kernel is obtained by completing the normal form (E.26g) to
coordinates in $\mathbb Z_p$. Conjugation (E.26e) extends
continuously: its formulas use addition, multiplication, and
$\binom a2$, the continuous integer-valued polynomial on
$\mathbb Z_p$. The finite cosets of $F$ give the corresponding
continuous extension over $S_3$.
The quotient (E.26c) receives (E.26a) by reducing its three
kernel coordinates modulo three. It therefore also receives
the profinite completion of (E.26a), continuously. A section
of that completion would give a section of (E.26c).

To construct (E.26c) completely, over $\mathbb F_3$ the law
(E.26g) is the Heisenberg group of order $27$ and exponent
three. It is exactly $F/(\gamma_3F F^3)$: collection bounds
its order by $27$, and the matrix model realizes all $27$
normal forms. The quotient has six such cosets, hence order
$162$. With states $x^a y^b c^k\rho^i\sigma^j$, right
multiplication by $\sigma$ toggles $j$. Right multiplication
by $\rho$ increments $i$ when $j=0$, and when $j=1$ it
multiplies the kernel coordinate by $\rho^i x\rho^{-i}$ and
replaces $i$ by $i+2$ modulo three. The three kernel factors are
$x,y,x^{-1}y^{-1}$. These rules construct its regular
permutation representation without a presentation-order guess.
The program \texttt{code/e\_splitting.py} checks closure of the
$162$ normal-form permutations under both generators, all
$4\,251\,528$ associativity triples, and all $27^2=729$
pairs of possible section-generator images. There are no
sections. Its abelian quotient of order $54$ has exactly nine
sections; all $81$ pairs lifting these have their
$\bar\rho$ image of order nine. The chosen complement has
preimage $\langle r',\sigma\rangle=D_{18}$, with
$(r')^3=c$. The code also checks all $216$ two-cocycle
identities and all $243$ normalized one-cochains modulo three:
none has coboundary $f$.

For $j\ge2$, the commutator identities
$[\gamma_aF,\gamma_bF]\subset\gamma_{a+b}F$ give
$F^{(j)}\subset\gamma_{2^j}F\subset\gamma_3F$. Consequently
the identity on generators gives the exact quotient map
\[
 P/F^{(j)}\twoheadrightarrow P/\gamma_3F,
 \quad\ker=\gamma_3F/F^{(j)}.                               \tag{E.26j}
\]
A section of the source extension would compose to a section
of (E.26a). For the profinite version, use the continuous map
to the finite group (E.26c); its kernel contains the closed
$j$th derived subgroup. The same composition gives the
contradiction. This establishes the transfer to the requested
third-derived quotient by an actual surjection.

Finally, the normalized lift calculation (E.26i) works with
central coordinates in $\mathbb Z_p$, $\mathbb Q$, or
$\mathbb Q_\ell$. It has the unique solution $a=-1/3$ exactly
in the coefficient cases specified in the statement. The
parameter $b$ remains arbitrary. Any kernel conjugation
preserving the normalized abelian complement has zero image
in $M$, by $H^0(S_3,M)=0$, and is therefore $c^u$.
It changes $b$ to $b+2u$. Thus the conjugacy classes are
$\mathbb Z_p/2\mathbb Z_p$ for $p\ne3$, and a singleton
over the indicated characteristic-zero fields. Formula
(E.26d) gives their specified representative. In the
cochain description, $-i/3$ cancels (E.26h), so the comparison
to rational coefficients explicitly kills this integral
obstruction.
The rational group here is the class-two relative unipotent
quotient over $S_3$: its six cosets are copies of the
Heisenberg group (E.26g) over the field. To verify the universal
comparison, a homomorphism from $F$ to a class-two unipotent
group, with images $X,Y$, extends uniquely by
$x^a y^b c^k\mapsto X^aY^b[X,Y]^k$. Rational powers are given
by the finite exponential and logarithm polynomials.
The finite-coset relations then give the unique extension
of the original homomorphism from $P$ over $S_3$.
The program also checks the normalized central equations
$1+3a=0$ and the conjugacy relation $b\sim b+2u$ modulo
$p^e$ for $p=2,3,5,7$ and $1\le e\le4$.

We also record precisely how the marked-depth construction
in E.2 compares. There is a surjective Lie homomorphism
\[
 L/D^3L\longrightarrow L/\gamma_3L,
 \qquad \ker=\gamma_3L/D^3L,
\]
because a monomial with at least three occurrences of $y$
has total degree at least three. The specific full
permutation action does not descend to that marked quotient:
the swap $x\leftrightarrow y$ sends
$[y,[y,[y,x]]]\in D^3L$ to
$[x,[x,[x,y]]]$, a nonzero depth-one element modulo $D^3L$.
Thus (E.26j) is an $S_3$-extension comparison for the derived
quotients, whereas the displayed marked-depth map has the
stated Lie source and target. The rational section (E.26d)
and the integral obstruction (E.26i) retain these exact
domains and coefficient rings.
\end{proof}
